\begin{proofbox}
\textbf{\textcolor{CProof}{证明对象}}
以下先证明焦点在 $x$ 轴时的公式，设
\[
C:\frac{x^2}{a^2}-\frac{y^2}{b^2}=1,\quad a,b>0,\quad P(x_0,y_0).
\]
若有两条不同的实切线 $PA,PB$，记切点 $A(x_1,y_1), B(x_2,y_2)$。最后给出两个实切点存在条件的高中代数证明。

\medskip
\begin{minipage}{\linewidth}
\textbf{\textcolor{CProof}{方法一：最低门槛——切线式 + 两点确定直线}}
\textit{适用信号：已学过双曲线上一点的切线方程；优先用于课堂首次讲解。}
\begin{center}
% Exact model: a=2, b=1, P=(0,2), A/B=(+/-sqrt(5),-1/2).
\begin{tikzpicture}[x=1.25cm,y=1.25cm,>=Stealth,font=\small]
\path[use as bounding box] (-4.6,-2) rectangle (4.6,3.1);
\node[font=\small\bfseries] at (0,2.85) {方法一：从两条切线确定两个共线点};
\draw[->,gray!55] (-4.2,0)--(4.2,0) node[right] {$x$};
\draw[->,gray!55] (0,-1.8)--(0,1.5);
\node[right,fill=white,inner sep=1pt] at (0,1) {$y$};
\node[below left,fill=white,inner sep=1pt] at (0,0) {$O$};
\draw[blue,thick,domain=-1.3:1.3,samples=100] plot ({2*cosh(\x)},{sinh(\x)});
\draw[blue,thick,domain=-1.3:1.3,samples=100] plot ({-2*cosh(\x)},{sinh(\x)});
\draw[red!75!black,thick] (-3.15,{2-sqrt(5)*3.15/2})--(0,2)--(3.15,{2-sqrt(5)*3.15/2});
\draw[teal!80!black,very thick] (-4,-.5)--(4,-.5);
\fill (0,2) circle (1.8pt) node[above=4pt] {$P(x_0,y_0)$};
\fill ({sqrt(5)},-.5) circle (1.8pt) node[above right=3pt,fill=white,inner sep=1.5pt] {$A(x_1,y_1)$};
\fill ({-sqrt(5)},-.5) circle (1.8pt) node[above left=3pt,fill=white,inner sep=1.5pt] {$B(x_2,y_2)$};
\node[red!75!black,fill=white,inner sep=2pt] at (-2,1.6) {$P\in PB$};
\node[red!75!black,fill=white,inner sep=2pt] at (2,1.6) {$P\in PA$};
\node[teal!80!black,below=4pt,fill=white,inner sep=2pt] at (0,-.5) {$A,B\in l$};
\node[blue,fill=white,inner sep=2pt] at (0,-1.6) {$C:\ x^2/4-y^2=1$};
\end{tikzpicture}

\end{center}
看图：$P$ 分别在两条红色切线上；代入两次，就把 $A,B$ 都放到绿色直线 $l$ 上。

\end{minipage}
\begin{enumerate}[leftmargin=*,itemsep=0.4em]
\item \textbf{在切点 $A$ 处写切线：}
\[PA:\quad \frac{x_1x}{a^2}-\frac{y_1y}{b^2}=1.\]
\textbf{为什么：}这是双曲线在曲线上一点 $A(x_1,y_1)$ 处的切线方程。
\item \textbf{用 $P$ 在这条切线上：}
\[\frac{x_1x_0}{a^2}-\frac{y_1y_0}{b^2}=1.\]
把乘法顺序交换，说明 $A(x_1,y_1)$ 落在直线
\[l:\quad \frac{x_0x}{a^2}-\frac{y_0y}{b^2}=1\]
上。
\item \textbf{对 $B$ 重复：}由 $P\in PB$ 得
\[\frac{x_2x_0}{a^2}-\frac{y_2y_0}{b^2}=1,\]
所以 $B$ 也在 $l$ 上。$A\ne B$，于是由两点确定一条直线，\textbf{$AB=l$}。证毕。
\end{enumerate}
\textbf{直观回扣：}由“$P$ 在切线 $PA$ 上”变为“$A$ 在 $l$ 上”，只用了一次代入和一次乘法交换。

\medskip
\begin{minipage}{\linewidth}
\textbf{\textcolor{CProof}{方法二：最快——先写候选直线，再反向验证}}
\textit{适用信号：已知 $P$，要求迅速写出并证明切点弦方程；已确认存在两个实切点。}
\begin{center}
% Exact model: a=2, b=1, P=(0,2), A/B=(+/-sqrt(5),-1/2).
\begin{tikzpicture}[x=1.25cm,y=1.25cm,>=Stealth,font=\small]
\path[use as bounding box] (-4.6,-2) rectangle (4.6,3.1);
\node[font=\small\bfseries] at (0,2.85) {方法二：从候选直线的交点反查切线};
\draw[->,gray!55] (-4.2,0)--(4.2,0) node[right] {$x$};
\draw[->,gray!55] (0,-1.8)--(0,1.5);
\node[right,fill=white,inner sep=1pt] at (0,1) {$y$};
\node[below left,fill=white,inner sep=1pt] at (0,0) {$O$};
\draw[blue,thick,domain=-1.3:1.3,samples=100] plot ({2*cosh(\x)},{sinh(\x)});
\draw[blue,thick,domain=-1.3:1.3,samples=100] plot ({-2*cosh(\x)},{sinh(\x)});
\draw[red!75!black,thick] (-3.15,{2-sqrt(5)*3.15/2})--(0,2)--(3.15,{2-sqrt(5)*3.15/2});
\draw[teal!80!black,very thick] (-4,-.5)--(4,-.5);
\fill (0,2) circle (1.8pt) node[above=4pt] {$P(x_0,y_0)$};
\fill ({sqrt(5)},-.5) circle (1.8pt) node[above right=3pt,fill=white,inner sep=1.5pt] {$T(u,v)$};
\fill ({-sqrt(5)},-.5) circle (1.8pt) node[above left=3pt,fill=white,inner sep=1.5pt] {$B$};
\node[red!75!black,fill=white,inner sep=2pt] at (-2,1.6) {$B$ 处切线过 $P$};
\node[red!75!black,fill=white,inner sep=2pt] at (2,1.6) {$T$ 处切线过 $P$};
\node[teal!80!black,below=4pt,fill=white,inner sep=2pt] at (0,-.5) {候选直线 $l$ 与 $C$ 交于 $T,B$};
\node[blue,fill=white,inner sep=2pt] at (0,-1.6) {$C:\ x^2/4-y^2=1$};
\end{tikzpicture}

\end{center}
看图：先取绿色直线上的交点 $T$，再检查 $T$ 处的红色切线是否经过 $P$。

\end{minipage}
\begin{enumerate}[leftmargin=*,itemsep=0.3em]
\item \textbf{先写：} $l:\ \dfrac{x_0x}{a^2}-\dfrac{y_0y}{b^2}=1$。
\item \textbf{反查切线：}任取 $l$ 与 $C$ 的交点 $T(u,v)$，则 $\dfrac{x_0u}{a^2}-\dfrac{y_0v}{b^2}=1$。而 $C$ 在 $T$ 处的切线是 $\dfrac{ux}{a^2}-\dfrac{vy}{b^2}=1$，代入 $P$ 恰好成立，故 $PT$ 是切线。
\item \textbf{收尾：}若 $l$ 与 $C$ 交于两个不同实点 $A,B$，它们正是从 $P$ 引出的两个切点，故 $l=AB$。交点确有两个的条件由后面的附证独立核验；反过来，任一切点也满足 $l$ 的方程，因此不会漏掉切点。证毕。
\end{enumerate}


\medskip
\begin{minipage}{\linewidth}
\textbf{\textcolor{CProof}{方法三：常规通用——设一般直线、用判别式证明}}
\textit{适用信号：学生尚未掌握切线式，或教师希望从高中二次方程知识完整推出结论；包含竖直切线。}
\begin{center}
% Exact model: P=(2,2); tangent points (2,0), (-10/3,-4/3).
\begin{tikzpicture}[x=1.25cm,y=1.25cm,>=Stealth,font=\small]
\path[use as bounding box] (-4.6,-2.7) rectangle (4.6,3.15);
\node[font=\small\bfseries] at (0,2.9) {方法三：一般直线式同时覆盖斜切线与竖直切线};
\draw[->,gray!55] (-4.2,0)--(4.2,0) node[right] {$x$};
\draw[->,gray!55] (0,-2)--(0,2.4) node[right] {$y$};
\node[below left] at (0,0) {$O$};
\draw[blue,thick,domain=-1.45:1.45,samples=100] plot ({2*cosh(\x)},{sinh(\x)});
\draw[blue,thick,domain=-1.45:1.45,samples=100] plot ({-2*cosh(\x)},{sinh(\x)});
\draw[teal!80!black,very thick] (-4,-1.5)--(3,.25);
\node[teal!80!black,below=5pt,fill=white,inner sep=2pt] at (-.8,-.7) {切点弦 $l$};
\draw[red!75!black,thick] (-4,-1.75)--(2,2);
\draw[red!75!black,thick,dashed] (2,-1.8)--(2,2.4);
\fill (2,2) circle (1.8pt) node[above right=3pt] {$P(x_0,y_0)$};
\fill (-10/3,-4/3) circle (1.8pt);
\node[fill=white,inner sep=2pt] at (-2.2,-1.85) {$T(a^2u,-b^2v)$};
\fill (2,0) circle (1.8pt);
\node[fill=white,inner sep=2pt] at (2.8,.35) {$(a,0)$};
\node[red!75!black,fill=white,inner sep=2pt] at (-1.4,1.35) {$L:\ ux+vy=1$};
\node[red!75!black,fill=white,inner sep=2pt] at (3.4,-.5) {$v=0:\ x=a$};
\node[blue,fill=white,inner sep=2pt] at (0,-2.45) {$C:\ x^2/4-y^2=1,\quad P=(2,2)$};
\end{tikzpicture}

\end{center}
看图：斜切线和竖直切线都能写成 $ux+vy=1$；它们的切点同在绿色直线上。

\end{minipage}
\begin{enumerate}[leftmargin=*,itemsep=0.42em]
\item \textbf{设一条切线：}双曲线的有限切线都不经过原点（过原点的直线与双曲线若相交，则交点成中心对称的一对，不可能只有一个切点）。因此可把切线化为
\[L:\quad ux+vy=1.\]
这里 $u\ne0$：若 $u=0$，则 $y$ 是常数，代入双曲线得到成对的 $x$ 值；没有恰为一个有限交点的相切情形。
\item \textbf{与双曲线联立：}由 $x=(1-vy)/u$，代入 $x^2/a^2-y^2/b^2=1$，整理得
\[
(b^2v^2-a^2u^2)y^2-2b^2vy+b^2(1-a^2u^2)=0.
\]
它是关于 $y$ 的方程。若二次项系数为零，方程只剩一次项，所得唯一交点不是相切的重合交点，须先排除。真正相切时二次项系数不为零，且两个交点重合，因此判别式为零：
\[
\begin{aligned}
\Delta&=(-2b^2v)^2-4(b^2v^2-a^2u^2)b^2(1-a^2u^2)\\
&=4a^2b^2u^2(1-a^2u^2+b^2v^2)=0.
\end{aligned}
\]
因 $a,b>0$ 且 $u\ne0$，可得
\[\boxed{a^2u^2-b^2v^2=1}.\]
此时二次项系数等于 $-1$，确实是有重根的二次方程，不会把一次方程误认为切线。
\item \textbf{求切点坐标：}二次方程的重根为
\[y_T=\frac{-(-2b^2v)}{2(-1)}=-b^2v.\]
再代回 $ux+vy=1$，得到
\[x_T=\frac{1+b^2v^2}{u}=a^2u.\]
所以切点 $T=(a^2u,-b^2v)$，也即 $u=x_T/a^2,\ v=-y_T/b^2$。
\item \textbf{利用切线经过 $P$：}因为 $ux_0+vy_0=1$，故
\[
\frac{x_0x_T}{a^2}-\frac{y_0y_T}{b^2}=ux_0+vy_0=1.
\]
于是任意切点 $T$ 都在直线 $l:\ \dfrac{x_0x}{a^2}-\dfrac{y_0y}{b^2}=1$ 上。对两个不同切点 $A,B$ 都成立，因此 $l=AB$。证毕。
\end{enumerate}
\textbf{补充：}当 $v=0$ 时，上面的推导仍成立，给出 $x=\pm a$ 的竖直切线；不必用斜率式 $y=kx+d$，因而不会漏掉这一特殊情形。

\medskip
\begin{minipage}{\linewidth}
\textbf{\textcolor{CProof}{附证：两个不同的有限实切点何时存在？}}
\begin{center}
\begin{tikzpicture}[font=\small,>=Stealth]
\begin{scope}[shift={(0,0)},x=.65cm,y=.65cm]
\node[font=\small\bfseries] at (0,3) {$P=(0,2),\ q=-4$};
\draw[->,gray!55] (-4.3,0)--(4.5,0);
\draw[->,gray!55] (0,-2)--(0,2.4);
\draw[blue,thick,domain=-1.4:1.4,samples=90] plot ({2*cosh(\x)},{sinh(\x)});
\draw[blue,thick,domain=-1.4:1.4,samples=90] plot ({-2*cosh(\x)},{sinh(\x)});
\draw[teal!80!black,very thick] (-4,-.5)--(4,-.5);
\draw[red!75!black,dashed,thick] (-3.1,{2-sqrt(5)*3.1/2})--(0,2)--(3.1,{2-sqrt(5)*3.1/2});
\fill (0,2) circle (2pt) node[above right=2pt,fill=white,inner sep=1pt] {$P$};
\fill ({sqrt(5)},-.5) circle (2.5pt);
\fill ({-sqrt(5)},-.5) circle (2.5pt);
\node at (0,-2.7) {两个切点};
\end{scope}
\begin{scope}[shift={(6.5,0)},x=.65cm,y=.65cm]
\node[font=\small\bfseries] at (0,3) {$P=(2,1),\ q=0$};
\draw[->,gray!55] (-4.3,0)--(4.5,0);
\draw[->,gray!55] (0,-2)--(0,2.4);
\draw[blue,thick,domain=-1.4:1.4,samples=90] plot ({2*cosh(\x)},{sinh(\x)});
\draw[blue,thick,domain=-1.4:1.4,samples=90] plot ({-2*cosh(\x)},{sinh(\x)});
\draw[teal!80!black,very thick] (-1,-1.5)--(4.5,1.25);
\draw[red!75!black,dashed,thick] (2,-1.7)--(2,2.1);
\fill (2,1) circle (2pt) node[above right=2pt,fill=white,inner sep=1pt] {$P$};
\fill (2,0) circle (2.5pt);
\node at (0,-2.7) {一个有限切点};
\end{scope}
\begin{scope}[shift={(0,-5.1)},x=.65cm,y=.65cm]
\node[font=\small\bfseries] at (0,3) {$P=(2,0),\ q=1$};
\draw[->,gray!55] (-4.3,0)--(4.5,0);
\draw[->,gray!55] (0,-2)--(0,2.4);
\draw[blue,thick,domain=-1.4:1.4,samples=90] plot ({2*cosh(\x)},{sinh(\x)});
\draw[blue,thick,domain=-1.4:1.4,samples=90] plot ({-2*cosh(\x)},{sinh(\x)});
\draw[teal!80!black,very thick] (2,-1.7)--(2,2.1);
\draw[red!75!black,dashed,thick] (2,-1.7)--(2,2.1);
\fill (2,0) circle (2pt) node[above right=2pt,fill=white,inner sep=1pt] {$P$};
\fill (2,0) circle (2.5pt);
\node at (0,-2.7) {一个切点（即 $P$）};
\end{scope}
\begin{scope}[shift={(6.5,-5.1)},x=.65cm,y=.65cm]
\node[font=\small\bfseries] at (0,3) {$P=(3,1),\ q=5/4$};
\draw[->,gray!55] (-4.3,0)--(4.5,0);
\draw[->,gray!55] (0,-2)--(0,2.4);
\draw[blue,thick,domain=-1.4:1.4,samples=90] plot ({2*cosh(\x)},{sinh(\x)});
\draw[blue,thick,domain=-1.4:1.4,samples=90] plot ({-2*cosh(\x)},{sinh(\x)});
\draw[teal!80!black,very thick] (-1,-1.75)--(4.5,2.375);
\fill (3,1) circle (2pt) node[below right=3pt,fill=white,inner sep=1pt] {$P$};
\node at (0,-2.7) {没有实切点};
\end{scope}
\node[align=center] at (3.25,-7.8) {同一双曲线 $x^2/4-y^2=1$，各图同尺度\\蓝：双曲线\quad 绿实线：代换直线\quad 红虚线：实切线};
\end{tikzpicture}

\end{center}
\end{minipage}
上图把候选直线的交点数直接呈现出来；下面用代数证明这些分界。
为看清条件，仅在这一小段令
\[
X=\frac{x}{a},\quad Y=\frac{y}{b},\quad U=\frac{x_0}{a},\quad V=\frac{y_0}{b},\quad q=U^2-V^2.
\]
双曲线变为 $X^2-Y^2=1$，候选直线变为 $UX-VY=1$。

\textbf{第1步（$q\ne0$）：}候选直线上所有点都可写为
\[
X=\frac{U}{q}+Vt,\qquad Y=\frac{V}{q}+Ut\qquad(t\in\mathbb R).
\]
因为 $UX-VY=(U^2-V^2)/q=1$，且方向 $(V,U)$ 沿着这条直线。代回双曲线，交叉项抵消：
\[
X^2-Y^2=\frac{1}{q}-qt^2=1
\quad\Longleftrightarrow\quad
\boxed{t^2=\frac{1-q}{q^2}}.
\]
因此：$q<1$ 时有两个不同实数 $t$（两个有限交点）；$q=1$ 时只有一个；$q>1$ 时无实数解。结合方法二，两个交点就对应两个切点。

\textbf{第2步（$q=0$）：}此时 $U=V$ 或 $U=-V$。若 $U=V=0$，直线方程为 $0=1$，不可能有切点。若 $U=V\ne0$，则
\[U(X-Y)=1,\quad (X-Y)(X+Y)=1,\]
从而 $X-Y=1/U,\ X+Y=U$，只确定\textbf{一个}有限交点；$U=-V\ne0$ 的情形相同（交换 $X-Y$ 与 $X+Y$）。故 $q=0$ 不能产生两个不同有限切点。

综上，\textbf{存在两个不同有限实切点的充要条件为 $q<1$ 且 $q\ne0$}。整个附证只用直线方程、平方差公式和一元二次方程的实数解。

\medskip
\textbf{\textcolor{CProof}{焦点在 $y$ 轴的推广}}
把以上推导中的 $x$ 与 $y$ 对调，即得
\[
\frac{y^2}{a^2}-\frac{x^2}{b^2}=1
\quad\Longrightarrow\quad
\boxed{AB:\ \frac{y_0y}{a^2}-\frac{x_0x}{b^2}=1},
\]
其中 $q=y_0^2/a^2-x_0^2/b^2$，仍须 $q<1$ 且 $q\ne0$。
\end{proofbox}
